\begin{trapbox}

	\begin{description}[style=nextline, leftmargin=2.8em, labelsep=0.5em, itemsep=0.35em]
		\item[\textbf{\textcolor{CTrap}{易错点一（混淆离心角极角）}}]
			错误认为参数方程中的 $\theta$ 就是 $OP$ 与 $x$ 轴的正向夹角。
		\item[\textbf{\textcolor{CTrap}{正确理解（区分两种角）：}}]
			离心角 $\theta$ 与极角 $\varphi$ 的关系是 $\tan\varphi=\frac{b}{a}\tan\theta$，二者一般不相等。
		\item[\textbf{\textcolor{CTrap}{错因分析（概念不清）：}}]
			将参数方程与极坐标方程混淆，忽视离心角来自辅助圆。
	\end{description}

	\medskip
	\begin{description}[style=nextline, leftmargin=2.8em, labelsep=0.5em, itemsep=0.35em]
		\item[\textbf{\textcolor{CTrap}{易错点二（忽视$\theta$范围）}}]
			求解具体问题时忽略 $\theta$ 的取值范围，导致多解或漏解。
		\item[\textbf{\textcolor{CTrap}{正确理解（根据象限定范围）：}}]
			通常 $\theta\in[0,2\pi)$，但需根据点所在象限选择适当的 $\theta$ 范围。
		\item[\textbf{\textcolor{CTrap}{错因分析（忽视周期性）：}}]
			参数 $\theta$ 与三角函数的周期相关，需结合题目限制确定范围。
	\end{description}

	\medskip
	\begin{description}[style=nextline, leftmargin=2.8em, labelsep=0.5em, itemsep=0.35em]
		\item[\textbf{\textcolor{CTrap}{易错点三（忽略有界性）}}]
			在求最值时直接令表达式等于某值，忽视 $\sin\theta,\cos\theta$ 有界。
		\item[\textbf{\textcolor{CTrap}{正确理解（利用有界性）：}}]
			使用 $|\sin\theta|\le1,|\cos\theta|\le1$ 作为最值依据。
		\item[\textbf{\textcolor{CTrap}{错因分析（思维定势）：}}]
			习惯用代数方法求解最值，忘记三角函数的有界约束。
	\end{description}

\end{trapbox}
